\chapter{矩阵}
\section{习题}
\begin{enumerate}
  \item 计算 $AB$，$AB-BA$．设
  \begin{enumerate}
    \item $A=\begin{pmatrix}3&1&1\\2&1&2\\1&2&3\end{pmatrix}$，$B=\begin{pmatrix}1&1&-1\\2&-1&0\\1&0&1\end{pmatrix}$；
  \anssub
    \item $A=\begin{pmatrix}a&b&c\\c&b&a\\1&1&1\end{pmatrix}$，$B=\begin{pmatrix}1&a&c\\1&b&b\\1&c&a\end{pmatrix}$，
  \anssub
  \end{enumerate}
  \item 计算：
  \begin{enumerate}
    \item $\begin{pmatrix}2&1&1\\3&1&0\\0&1&2\end{pmatrix}^2$；
  \anssub
    \item $\begin{pmatrix}3&2\\-4&-2\end{pmatrix}^5$；
  \anssub
    \item $\begin{pmatrix}1&1\\0&1\end{pmatrix}^n$；
  \anssub
    \item $\begin{pmatrix}\cos\varphi&-\sin\varphi\\\sin\varphi&\cos\varphi\end{pmatrix}^n$；
  \anssub
    \item $(2,3,-1)\begin{pmatrix}1\\-1\\-1\end{pmatrix}$，$\begin{pmatrix}1\\-1\\-1\end{pmatrix}(2,3,-1)$；
  \anssub
    \item $(x,y,1)\begin{pmatrix}a_{11}&a_{12}&b_1\\a_{12}&a_{22}&b_2\\b_1&b_2&c\end{pmatrix}\begin{pmatrix}x\\y\\1\end{pmatrix}$；
  \anssub
    \item $\begin{pmatrix}1&-1&-1&-1\\-1&1&-1&-1\\-1&-1&1&-1\\-1&-1&-1&1\end{pmatrix}^2$，$\begin{pmatrix}1&-1&-1&-1\\-1&1&-1&-1\\-1&-1&1&-1\\-1&-1&-1&1\end{pmatrix}^n$；
  \anssub
    \item $\begin{pmatrix}\lambda&1&0\\0&\lambda&1\\0&0&\lambda\end{pmatrix}^n$．
  \anssub
  \end{enumerate}
  \item 若 $f(\lambda)=a_0\lambda^m+a_1\lambda^{m-1}+\cdots+a_m$，$A$ 是一个 $n\times n$ 矩阵，定义 $f(A)=a_0A^m+a_1A^{m-1}+\cdots+a_mE$．试求 $f(A)$．设
  \begin{enumerate}
    \item $f(\lambda)=\lambda^2-\lambda-1$，$A=\begin{pmatrix}2&1&1\\3&1&2\\1&-1&0\end{pmatrix}$；
  \anssub
    \item $f(\lambda)=\lambda^2-5\lambda+3$，$A=\begin{pmatrix}2&-1\\-3&3\end{pmatrix}$．
  \anssub
  \end{enumerate}
  \item 如果 $AB=BA$，矩阵 $B$ 就称为与矩阵 $A$ 可交换．分别求所有与 $A$ 可交换的矩阵．设
  \begin{enumerate}
    \item $A=\begin{pmatrix}1&1\\0&1\end{pmatrix}$；
  \anssub
    \item $A=\begin{pmatrix}1&0&0\\0&1&2\\3&1&2\end{pmatrix}$；
  \anssub
    \item $A=\begin{pmatrix}0&1&0\\0&0&1\\0&0&0\end{pmatrix}$，
  \anssub
  \end{enumerate}
  \item 设
  \[
    A=\begin{pmatrix}
      a_1&0&\cdots&0\\
      0&a_2&\cdots&0\\
      \vdots&\vdots&&\vdots\\
      0&0&\cdots&a_n
    \end{pmatrix},
  \]
  其中 $a_i\neq a_j(i\neq j,i,j=1,2,\cdots,n)$．证明：与 $A$ 可交换的矩阵只能是对角矩阵．
  \ansproof
  \item 设
  \[
    A=\begin{pmatrix}
      a_1E_1&&&O\\
      &a_2E_2&&\\
      &&\ddots&\\
      O&&&a_rE_r
    \end{pmatrix},
  \]
  其中 $a_i\neq a_j(i\neq j,i,j=1,2,\cdots,r)$，$E_i$ 是 $n_i$ 阶单位矩阵，$\sum\limits_{i=1}^{r}n_i=n$．证明：与 $A$ 可交换的矩阵只能是准对角矩阵
  \[
    \begin{pmatrix}
      A_1&&&O\\
      &A_2&&\\
      &&\ddots&\\
      O&&&A_r
    \end{pmatrix},
  \]
  其中 $A_i,i=1,2,\cdots,r$，是 $n_i$ 阶矩阵．
  \ansproof
  \item 用 $E_{ij}$ 表示第 $i$ 行第 $j$ 列的元素为 $1$，而其余元素全为零的 $n\times n$ 矩阵，而 $A=(a_{ij})_{n\times n}$．证明：
  \begin{enumerate}
    \item 如果 $AE_{12}=E_{12}A$，那么当 $k\neq1$ 时 $a_{k1}=0$，当 $k\neq2$ 时 $a_{2k}=0$；
  \anssub
    \item 如果 $AE_{ij}=E_{ij}A$，那么当 $k\neq i$ 时 $a_{ki}=0$，当 $k\neq j$ 时 $a_{jk}=0$，且 $a_{ii}=a_{jj}$；
  \anssub
    \item 如果 $A$ 与所有的 $n$ 阶矩阵可交换，那么 $A$ 一定是数量矩阵，即 $A=aE$．
  \anssub
  \end{enumerate}
  \item 如果 $AB=BA$，$AC=CA$，证明：$A(B+C)=(B+C)A$，$A(BC)=(BC)A$．
  \ansproof
  \item 如果 $A=\dfrac{1}{2}(B+E)$，证明：$A^2=A$ 当且仅当 $B^2=E$．
  \ansproof
  \item 矩阵 $A$ 称为对称的，如果 $A^{\mathrm{T}}=A$．证明：如果 $A$ 是实对称矩阵且 $A^2=O$，那么 $A=O$．
  \ansproof
  \item 设 $A,B$ 都是 $n\times n$ 对称矩阵．证明：$AB$ 也对称当且仅当 $A,B$ 可交换．
  \ansproof
  \item 矩阵 $A$ 称为反称的，如果 $A^{\mathrm{T}}=-A$．证明：任一 $n\times n$ 矩阵都可表为一对称矩阵与一反称矩阵之和．
  \ansproof
  \item 设 $s_k=x_1^k+x_2^k+\cdots+x_n^k$，$k=0,1,2,\cdots$，$a_{ij}=s_{i+j-2}$，$i,j=1,2,\cdots,n$．证明：行列式
  \[
    |a_{ij}|=\prod_{i<j}(x_i-x_j)^2.
  \]
  \ansproof
  \item 设 $A$ 是 $n\times n$ 矩阵．证明：存在一个 $n\times n$ 非零矩阵 $B$，使 $AB=O$ 的充要条件是 $|A|=0$．
  \ansproof
  \item 设 $A$ 是 $n\times n$ 矩阵．证明：如果对任一 $n$ 维向量 $X=\begin{pmatrix}x_1\\x_2\\\vdots\\x_n\end{pmatrix}$ 都有 $AX=0$，那么 $A=O$．
  \ansproof
  \item 设 $B$ 为一 $r\times r$ 矩阵，$C$ 为一 $r\times n$ 矩阵，且秩$(C)=r$．证明：
  \begin{enumerate}
    \item 如果 $BC=O$，那么 $B=O$；
  \anssub
    \item 如果 $BC=C$，那么 $B=E$．
  \anssub
  \end{enumerate}
  \item 证明：
  \[
    \text{秩}(A+B)\leqslant\text{秩}(A)+\text{秩}(B).
  \]
  \ansproof
  \item 设 $A,B$ 为 $n\times n$ 矩阵．证明：如果 $AB=O$，那么
  \[
    \text{秩}(A)+\text{秩}(B)\leqslant n.
  \]
  \ansproof
  \item 证明：如果 $A^k=O$，那么
  \[
    (E-A)^{-1}=E+A+A^2+\cdots+A^{k-1}.
  \]
  \ansproof
  \item 求下列 $A^{-1}$：
  \begin{enumerate}
    \item $A=\begin{pmatrix}a&b\\c&d\end{pmatrix}$，$ad-bc=1$；
  \anssub
    \item $A=\begin{pmatrix}1&1&-1\\2&1&0\\1&-1&0\end{pmatrix}$；
  \anssub
    \item $A=\begin{pmatrix}2&2&3\\1&-1&0\\-1&2&1\end{pmatrix}$；
  \anssub
    \item $A=\begin{pmatrix}1&2&3&4\\2&3&1&2\\1&1&1&-1\\1&0&-2&-6\end{pmatrix}$；
  \anssub
    \item $A=\begin{pmatrix}1&1&1&1\\1&1&-1&-1\\1&-1&1&-1\\1&-1&-1&1\end{pmatrix}$；
  \anssub
    \item $A=\begin{pmatrix}3&3&-4&-3\\0&6&1&1\\5&4&2&1\\2&3&3&2\end{pmatrix}$；
  \anssub
    \item $A=\begin{pmatrix}1&3&-5&7\\0&1&2&-3\\0&0&1&2\\0&0&0&1\end{pmatrix}$；
  \anssub
    \item $A=\begin{pmatrix}2&1&0&0\\3&2&0&0\\5&7&1&8\\-1&-3&-1&-6\end{pmatrix}$；
  \anssub
    \item $A=\begin{pmatrix}0&0&1&-1\\0&3&1&4\\2&7&6&-1\\1&2&2&-1\end{pmatrix}$；
  \anssub
    \item $A=\begin{pmatrix}2&1&0&0&0\\0&2&1&0&0\\0&0&2&1&0\\0&0&0&2&1\\0&0&0&0&2\end{pmatrix}$．
  \anssub
  \end{enumerate}
  \item 设
  \[
    X=\begin{pmatrix}O&A\\C&O\end{pmatrix},
  \]
  已知 $A^{-1}$，$C^{-1}$ 存在，求 $X^{-1}$．
  \anscompute
  \item 设
  \[
    X=\begin{pmatrix}
      0&a_1&0&\cdots&0&0\\
      0&0&a_2&\cdots&0&0\\
      \vdots&\vdots&\vdots&&\vdots&\vdots\\
      0&0&0&\cdots&0&a_{n-1}\\
      a_n&0&0&\cdots&0&0
    \end{pmatrix},
  \]
  其中 $a_i\neq0(i=1,2,\cdots,n)$，求 $X^{-1}$．
  \anscompute
  \item 求下列矩阵 $X$：
  \begin{enumerate}
    \item $\begin{pmatrix}2&5\\1&3\end{pmatrix}X=\begin{pmatrix}4&-6\\2&1\end{pmatrix}$；
  \anssub
    \item $\begin{pmatrix}1&1&-1\\0&2&2\\1&-1&0\end{pmatrix}X=\begin{pmatrix}1&-1&1\\1&1&0\\2&1&1\end{pmatrix}$；
  \anssub
    \item
    \[
      \begin{pmatrix}
        1&1&1&\cdots&1&1\\
        0&1&1&\cdots&1&1\\
        0&0&1&\cdots&1&1\\
        \vdots&\vdots&\vdots&&\vdots&\vdots\\
        0&0&0&\cdots&0&1
      \end{pmatrix}_{n\times n}
      X=
      \begin{pmatrix}
        2&1&0&\cdots&0&0\\
        1&2&1&\cdots&0&0\\
        0&1&2&\cdots&0&0\\
        \vdots&\vdots&\vdots&&\vdots&\vdots\\
        0&0&0&\cdots&1&2
      \end{pmatrix}_{n\times n};
    \]
  \anssub
    \item $X\begin{pmatrix}1&1&-1\\0&2&2\\1&-1&0\end{pmatrix}=\begin{pmatrix}1&-1&1\\1&1&0\\2&1&1\end{pmatrix}$．
  \anssub
  \end{enumerate}
  \item 证明：
  \begin{enumerate}
    \item 如果 $A$ 可逆对称（反称），那么 $A^{-1}$ 也对称（反称）；
  \anssub
    \item 不存在奇数阶的可逆反称矩阵．
  \anssub
  \end{enumerate}
  \item 矩阵 $A=(a_{ij})$ 称为上（下）三角形矩阵，如果当 $i>j$（$i<j$）时有 $a_{ij}=0$．证明：
  \begin{enumerate}
    \item 两个上（下）三角形矩阵的乘积仍是上（下）三角形矩阵；
  \anssub
    \item 可逆的上（下）三角形矩阵的逆仍是上（下）三角形矩阵．
  \anssub
  \end{enumerate}
  \item 证明：
  \[
    |A^*|=|A|^{n-1},
  \]
  其中 $A$ 是 $n\times n$ 矩阵（$n\geqslant2$）．
  \ansproof
  \item 证明：如果 $A$ 是 $n\times n$（$n\geqslant2$）矩阵，那么
  \[
    \text{秩}(A^*)=
    \begin{cases}
      n,&\text{秩}(A)=n,\\
      1,&\text{秩}(A)=n-1,\\
      0,&\text{秩}(A)<n-1.
    \end{cases}
  \]
  \ansproof
  \item 用下列两种方法求矩阵
  \[
    A=\begin{pmatrix}
      1&1&1&1\\
      1&-1&1&-1\\
      1&1&-1&-1\\
      1&-1&-1&1
    \end{pmatrix}
  \]
  的逆矩阵：
  \begin{enumerate}
    \item 初等变换；
  \anssub
    \item 按 $A$ 中的划分，利用分块乘法的初等变换（注意各小块矩阵的特点）．
  \anssub
  \end{enumerate}
  \item 设 $A,B$ 分别是 $n\times m$ 和 $m\times n$ 矩阵．证明：
  \[
    \begin{vmatrix}E_m&B\\A&E_n\end{vmatrix}=|E_n-AB|=|E_m-BA|.
  \]
  \ansproof
  \item 设 $A,B$ 如上题，$\lambda\neq0$．证明：
  \[
    |\lambda E_n-AB|=\lambda^{n-m}|\lambda E_m-BA|.
  \]
  \ansproof
  \item 设数域 $P$ 上 $n$ 阶矩阵 $C=(\varepsilon_n,\varepsilon_1,\varepsilon_2,\cdots,\varepsilon_{n-1})$，其中 $\varepsilon_j(j=1,2,\cdots,n)$ 表示第 $j$ 个分量为 $1$，其余分量为 $0$ 的 $n$ 维列向量，求 $C^m$，其中 $m$ 为任意正整数．
  \anscompute
\end{enumerate}

\section{补充题}
\begin{enumerate}
  \item 设 $A$ 是 $n\times n$ 矩阵，秩$(A)=1$．证明：
  \begin{enumerate}
    \item $A=\begin{pmatrix}a_1\\a_2\\\vdots\\a_n\end{pmatrix}(b_1,b_2,\cdots,b_n)$；
  \anssub
    \item $A^2=kA$．
  \anssub
  \end{enumerate}
  \item 设 $A$ 是 $2\times2$ 矩阵．证明：如果 $A^l=O$，$l\geqslant2$，那么 $A^2=O$．
  \ansproof
  \item 设 $A$ 是 $n\times n$ 矩阵．证明：如果 $A^2=E$，那么
  \[
    \text{秩}(A+E)+\text{秩}(A-E)=n.
  \]
  \ansproof
  \item 设 $A$ 是 $n\times n$ 矩阵，且 $A^2=A$．证明：
  \[
    \text{秩}(A)+\text{秩}(A-E)=n.
  \]
  \ansproof
  \item 证明：
  \[
    (A^*)^*=|A|^{n-2}A,
  \]
  其中 $A$ 是 $n\times n$（$n>2$）矩阵．
  \ansproof
  \item 设 $A,B,C,D$ 都是 $n\times n$ 矩阵，且 $|A|\neq0$，$AC=CA$．证明：
  \[
    \begin{vmatrix}A&B\\C&D\end{vmatrix}=|AD-CB|.
  \]
  \ansproof
  \item 设 $A$ 是一 $n\times n$ 矩阵，且秩$(A)=r$．证明：存在一 $n\times n$ 可逆矩阵 $P$，使 $PAP^{-1}$ 的后 $n-r$ 行全为零．
  \ansproof
  \item
  \begin{enumerate}
    \item 把矩阵
    \[
      \begin{pmatrix}a&0\\0&a^{-1}\end{pmatrix}
    \]
    表成形如
    \[
      \begin{pmatrix}1&x\\0&1\end{pmatrix}
      \quad\text{与}\quad
      \begin{pmatrix}1&0\\x&1\end{pmatrix}
      \tag{1}
    \]
    的矩阵的乘积．
  \anssub
    \item 设
    \[
      A=\begin{pmatrix}a&b\\c&d\end{pmatrix}
    \]
    为一复数矩阵，$|A|=1$．证明：$A$ 可以表成形如（1）式的矩阵的乘积．
  \anssub
  \end{enumerate}
  \item 设 $A$ 是一 $n\times n$ 矩阵，$|A|=1$．证明：$A$ 可以表成 $P(i,j(k))$ 这一类初等矩阵的乘积．
  \ansproof
  \item 设 $A=(a_{ij})_{s\times n}$，$B=(b_{ij})_{n\times m}$．证明：
  \[
    \text{秩}(AB)\geqslant\text{秩}(A)+\text{秩}(B)-n.
  \]
  \ansproof
  \item 矩阵的列（行）向量组如果是线性无关的，就称该矩阵为列（行）满秩的．证明：设 $A$ 是 $m\times r$ 矩阵，则 $A$ 是列满秩的充要条件为存在 $m\times m$ 可逆矩阵 $P$，使
  \[
    A=P\begin{pmatrix}E_r\\O\end{pmatrix}.
  \]
  同样地，$A$ 为行满秩的充要条件为存在 $r\times r$ 可逆矩阵 $Q$，使
  \[
    A=(E_m\ \ O)Q.
  \]
  \ansproof
  \item 证明：设 $m\times n$ 矩阵 $A$ 的秩为 $r$，则有 $m\times r$ 列满秩矩阵 $P$ 和 $r\times n$ 行满秩矩阵 $Q$，使 $A=PQ$．
  \ansproof
  \item 定义数域 $P$ 上的 $n$ 阶方阵 $X_{ij}(\alpha)=E+\alpha E_{ij}$，其中 $\alpha\in P$，$E$ 为单位矩阵，$E_{ij}$ 为 $(i,j)$ 元素为 $1$，其他元素为 $0$ 的方阵．证明：当 $n\geqslant2$ 时，对于任意 $\alpha\in P$，$1\leqslant i\neq j\leqslant n$，总存在 $n$ 阶可逆矩阵 $R,T$，使得 $X_{ij}(\alpha)=R^{-1}T^{-1}RT$．
  \ansproof
  \item
  \begin{enumerate}
    \item 证明：对于数域 $P$ 上任意 $n+1$ 个 $n$ 阶方阵 $A_1,A_2,\cdots,A_{n+1}$，总能找到不全为零的数 $k_1,k_2,\cdots,k_{n+1}\in P$，使得矩阵 $k_1A_1+k_2A_2+\cdots+k_{n+1}A_{n+1}$ 的秩小于 $n$；
  \anssub
    \item 对于数域 $P$ 上任意 $n(n>1)$ 个 $n$ 阶方阵 $A_1,A_2,\cdots,A_n$，是否总有不全为零的数 $k_1,k_2,\cdots,k_n\in P$，使得矩阵 $k_1A_1+k_2A_2+\cdots+k_nA_n$ 的秩小于 $n$？
  \anssub
  \end{enumerate}
  \item 设矩阵
  \[
    A(x,a,n)=\begin{pmatrix}
      x&a&a&\cdots&a\\
      a&x&a&\cdots&a\\
      a&a&x&\cdots&a\\
      \vdots&\vdots&\vdots&&\vdots\\
      a&a&a&\cdots&x
    \end{pmatrix}_{n\times n},
  \]
  于是有等式 $A(x,a,n)A(y,b,n)=A(z,c,n)$ 成立．
  \begin{enumerate}
    \item 试用 $x,y,a,b$ 表示 $z,c$；
  \anssub
    \item 判断矩阵 $A(0,1,4)$ 是否可逆，若可逆求出其逆矩阵．
  \anssub
  \end{enumerate}
  \item 设矩阵 $A=(a_{ij})$ 和 $Q$ 均为 $n$ 阶方阵，矩阵 $Q$ 为若干 $P(i,j)$ 这类初等矩阵的乘积，令 $B=QAQ^{\mathrm{T}}$．判断：$a_{ij}$ 在 $A$ 中的代数余子式 $A_{ij}$ 是否等于 $a_{ij}$ 在 $B$ 中的代数余子式？若相等，给出证明；若不相等，举出反例．
  \ansproof
  \item 设 $B$ 为数域 $P$ 上所有 $n$ 阶可逆上三角形矩阵的集合．对于 $1\leqslant i\leqslant n-1$，设 $N_i=P(i,i+1)$．令 $B_i=\{N_iMN_i\mid M\in B\}$．证明：集合 $B$ 不等于任意一个集合 $B_i(1\leqslant i\leqslant n-1)$．
  \ansproof
  \item 设 $A$ 为 $n$ 阶方阵，$A$ 的每一行每一列恰好有一个元素为 $1$ 或者为 $-1$，其余元素均为 $0$．证明：必存在正整数 $k$，使得 $A^k=E$．
  \ansproof
  \item 设 $M$ 为 $n$ 阶可逆矩阵．证明：存在一个 $n$ 阶可逆上三角形矩阵 $B$，使得矩阵 $BM$ 具有如下性质：对于任意 $1\leqslant i\leqslant n$，都存在 $BM$ 的一行，该行前 $i-1$ 个位置为 $0$ 且第 $i$ 个位置不为 $0$．
  \ansproof
  \item 设整数 $n\geqslant4$，$A$ 是一个数域 $P$ 上的 $n\times(n+1)$ 矩阵，$A_j$ 是去掉 $A$ 中第 $j$ 列的 $n$ 阶方阵，$|A_j|$ 表示 $A_j(1\leqslant j\leqslant n+1)$ 的行列式，$|A_3||A_4|\neq0$．证明：
  \[
    \frac{|A_1|^3|A_2|}{|A_3|^2|A_4|^2}
  \]
  在矩阵 $A$ 的初等行变换下保持不变．
  \ansproof
\end{enumerate}
